\documentclass[a4paper,12pt]{article}

\usepackage[french]{babel}
\usepackage[T1]{fontenc} 
\usepackage[latin1]{inputenc}
%\usepackage{amsmath}

\author{R�mi Synave, Stefka Gueorguieva, Pascal Desbarats}
\title{Manuel de la biblioth�que MATRIX}
\date{}

\begin{document}
\maketitle

\section{Utilisation}
Cette biblioth�que permet la manipulation de matrice et d'effectuer diff�rents calculs comme le calcul du d�terminant, l'addition, la soustraction ou la multiplication de deux matrices.\\
Cette biblioth�que utilise les matrices d�finies dans la biblioth�que GSL (http://www.gnu.org/software/gsl/).\\
Il faudra donc installer GSL avant d'utiliser cette biblioth�que.

\section{Fonctions}

\textbullet \texttt{double matrix\_determinant(gsl\_matrix *m)}\\
Calcul du d�terminant d'une matrice.\\
~\\
\underline{Param�tres et type de retour :}\\
\texttt{m} : matrice dont on cherche le d�terminant.\\
\texttt{retour} : d�terminant de la matrice.\\


\textbullet \texttt{double matrix\_frobenius\_norm(gsl\_matrix *m)}\\
Calcul de la norme de frobenius (aussi appel�e norme euclidienne) d'une matrice. La norme de frobenius d'une matrice $A$ est not�e : $||A||_F\\$
\textbf{Formule :} (http://mathworld.wolfram.com/FrobeniusNorm.html)\\
La norme de frobenius est la racine carr�e de la somme des �l�ments au carr�s.\\
Calcul de la norme de frobenius d'un matrice $A$ de $m$ lignes et $n$ colonnes.\\
$$||A||_F=\sqrt{\sum_{i=0}^m \sum_{j=0}^n a_{ij}^2}$$
~\\
\underline{Param�tres et type de retour :}\\
\texttt{m} : matrice dont on cherche la norme de Frobenius.\\
\texttt{retour} : norme de Frobenius.\\



\textbullet \texttt{gsl\_matrix* matrix\_add(gsl\_matrix *A, gsl\_matrix *B)}\\
Addition de deux matrices. Attention aux tailles des matrices qui doivent �tre compatibles avec le calcul.\\
~\\
\underline{Param�tres et type de retour :}\\
\texttt{A} : premi�re matrice.\\
\texttt{B} : seconde matrice.\\
\texttt{retour} : Nouvelle matrice contenant l'addition de A et B.\\


\textbullet \texttt{gsl\_matrix* matrix\_sub(gsl\_matrix *A, gsl\_matrix *B)}\\
Soustraction de deux matrices. Attention aux tailles des matrices qui doivent �tre compatibles avec le calcul.\\
~\\
\underline{Param�tres et type de retour :}\\
\texttt{A} : premi�re matrice.\\
\texttt{B} : seconde matrice.\\
\texttt{retour} : Nouvelle matrice contenant la soustraction de A et B.\\


\textbullet \texttt{gsl\_matrix* matrix\_mul\_scale(gsl\_matrix *A, double n)}\\
Multiplication d'une matrice par un scalaire.\\
~\\
\underline{Param�tres et type de retour :}\\
\texttt{A} : la matrice � multiplier.\\
\texttt{n} : le scalaire.\\
\texttt{retour} : Nouvelle matrice contenant la matrice multipli� par le scalaire.\\


\textbullet \texttt{gsl\_matrix* matrix\_mul(gsl\_matrix *A, gsl\_matrix *B)}\\
Multiplication de deux matrices. Attention aux tailles des matrices qui doivent �tre compatibles avec le calcul.\\
~\\
\underline{Param�tres et type de retour :}\\
\texttt{A} : premi�re matrice.\\
\texttt{B} : seconde matrice.\\
\texttt{retour} : Nouvelle matrice contenant la multiplication de A et B.\\


\textbullet \texttt{gsl\_matrix* matrix\_inverse(gsl\_matrix *m)}\\
Calcul de l'inverse d'une matrice.\\
~\\
\underline{Param�tres et type de retour :}\\
\texttt{m} : la matrice.\\
\texttt{retour} : Nouvelle matrice contenant l'inverse de m.\\


\textbullet \texttt{void matrix\_display(gsl\_matrix *m)}\\
Affichage d'une matrice.\\
~\\
\underline{Param�tres et type de retour :}\\
\texttt{m} : la matrice.\\
\texttt{retour} : aucun.\\


\textbullet \texttt{gsl\_matrix* matrix\_init(double *data, int nbline, int nbcol)}\\
Initialisation d'une matrice � l'aide d'un tableau de r��ls.\\
\textbf{Exemple d'utilisation :}
\begin{verbatim}
data = {1 2 3 4 5 6}
nbline = 3
nbcol = 2

          (1 2)
matrice = (3 4)
          (5 6)
\end{verbatim}
~\\
\underline{Param�tres et type de retour :}\\
\texttt{data} : donn�es pour remplir la matrice.\\
\texttt{nbline} : Nombre de lignes.\\
\texttt{nbcol} : Nombre de colonnes\\
\texttt{retour} : Nouvelle matrice contenant les donn�es du tableau data dans \texttt{nbline} lignes et \texttt{nbcol} colonnes.\\


\textbullet \texttt{gsl\_matrix* matrix\_transpose(gsl\_matrix *m)}\\
Calcul de la transpos� d'une matrice.\\
~\\
\underline{Param�tres et type de retour :}\\
\texttt{m} : la matrice.\\
\texttt{retour} : Nouvelle matrice contenant la transpose de m.\\


\end{document}
